\documentclass[11pt]{article}
\usepackage[margin=1in]{geometry}
\usepackage{amsmath,amssymb,amsthm,booktabs}
\usepackage[hidelinks]{hyperref}
\usepackage{xurl}
\setlength{\emergencystretch}{2em}
\newtheorem{theorem}{Theorem}
\newtheorem{lemma}[theorem]{Lemma}
\newtheorem{corollary}[theorem]{Corollary}
\newcommand{\rk}{\operatorname{rank}}
\newcommand{\mr}{\operatorname{mr}}
\title{Signed generic rank-one obstructions for sixteen knots}
\author{Written proof and exact Laurent controls}
\date{17 November 2026}
\begin{document}\maketitle

The generic rank-one obstruction for $13a_{373}$ extends to fifteen
further alternating knots. The argument uses signed Wirtinger
relations and independent root and covector families, without a
preserved form. It works over every field and every nonzero meridian
determinant, including unipotent meridians. Its augmentation consequence
concerns all meridian parameters admitted by the augmentation--KCH
correspondence, rather than only the involutory value $-1$.

Let
\[
\begin{aligned}
\mathcal K_{16}=\{&13a_{373},13a_{375},13a_{391},13a_{392},\\
&13a_{3264},13a_{3266},13a_{3275},13a_{3276},\\
&13a_{304},13a_{362},13a_{623},13a_{629},\\
&13a_{3232},13a_{3259},13a_{3374},13a_{3375}\}.
\end{aligned}
\]
The companion \path{GENERIC_KCH_OBSTRUCTION.tex} proves the first
case and supplies the arbitrary-pairing three-direction contradiction
used below. Its exact integer Laurent controls are an explicit
mathematical dependency. The present upper theorem is independent of
any positive augmentation witness.

\begin{theorem}\label{thm:all16}
For every $K\in\mathcal K_{16}$, every irreducible finite-dimensional
representation of its knot group with meridian displacement rank one
has dimension at most three, over every field.
No symmetric or Hermitian form is assumed.
\end{theorem}

\section{Signed pairs and independence}

Write a meridian image as
\[
 M_i=I-u_if_i,\qquad f_i(u_i)=d=1-\mu,
 \qquad M_i^{-1}=I+\mu^{-1}u_if_i.
\]
Here $u_i,f_i$ are nonzero and $\mu=\det M_i$ is a nonzero scalar,
the same for all conjugate meridians. The case $\mu=1$ is allowed.
At a signed crossing $(o,a,b,\varepsilon)$ one has
\[
 M_b=M_o^\varepsilon M_aM_o^{-\varepsilon},\qquad
 (u_b,f_b)=(M_o^\varepsilon u_a,f_aM_o^{-\varepsilon}).
\]
Equality of nonzero rank-one displacements changes a pair only by
$(u,f)\mapsto(cu,c^{-1}f)$ for a nonzero scalar $c$.

The four retained meridian seeds generate each knot group by a complete
literal diagram trace. Their root span is nonzero and invariant, and
their covectors' common kernel is fixed pointwise. Hence an irreducible
representation has dimension at most four. In dimension four the four
seed roots and the four seed covectors are both bases. We assume this
case and derive a contradiction.

All changes of seed systems used below are compositions of conjugating
one seed pair by an element whose root and covector belong to the spans
of the other three retained pairs. The replacement has coefficient one
on its old root and coefficient one on its old covector. Thus both basis
determinants are one. The exact checker independently computes these
two determinants in twelve unrestricted directed pairings and $\mu$,
using the actual signed crossing words for every diagram.
Only powers of $\mu$ occur as algebraic denominators.

\section{The four three-direction contradictions}

Identify the paired meridian matrices on arcs $0,4$ and on arcs
$6,7$. Choose their pairs to be $(U,F),(-U,-F)$ and $(V,G),(V,G)$.
Let $(W,J)$ be a third pair, on arcs $8,9,9,12$ in the four cases
$373,375,391,392$, respectively. Each corresponding evaluation is $d$.
Allow all six directed pairings
\[
 x=F(V),\quad X=G(U),\quad y=F(W),\quad Y=J(U),\quad
 z=G(W),\quad Z=J(V)
\]
to be unrestricted.

\begin{lemma}\label{lem:cores}
In each of these four signed reduced eight-crossing systems, the full
root span and the full covector span cannot both have dimension three.
\end{lemma}
\begin{proof}
The three selected pairs propagate every reduced pair, so each span
has dimension at most three. If both have dimension three, both selected
triples are independent. For $373$ the companion generic proof applies
directly: its contradiction uses only the reduced signed crossings and
the six free directed pairings, not the discarded four-seed clasp.
We prove the other three cases here. In the notation $M(T,L)$ below,
conjugation acts on a root and covector simultaneously and contragrediently.

For $375$ put
\[
 (T,L)=(M_V^{-1}W,JM_V),\qquad
 (A,H)=(M_UT,LM_U^{-1}),\qquad t=F(T),\quad\tau=L(U).
\]
Then $A=T-tU$, $H=L+(\tau/\mu)F$.
The signed propagation steps are
\[
 (6,9,10,-),(10,7,8,-),(4,10,11,+),
 (11,4,5,+),(8,11,12,-).
\]
They give the pairs $(T,L)$ on arc $10$, $(A,H)$ on arc $11$,
$M_8=M_T^{-1}M_VM_T$, and
\[
\begin{split}
 u_5&=(-1-t\tau/\mu)U+(\tau/\mu)T,\\
 f_5&=(-1-t\tau/\mu)F-tL.
\end{split}
\]
The coefficients of $u_{12},f_{12}$ on $U,F$ are $-t,\tau/\mu$,
since the arc $8$ pair uses only the other two coordinates.
The crossing $(9,12,0,+)$ says
\[
 u_{12}=\gamma M_W^{-1}U,\qquad
 f_{12}=\gamma^{-1}FM_W
\]
for $\gamma\ne0$. Thus $t=-\gamma$ and
$\tau=\mu\gamma^{-1}$, so $t\tau=-\mu$.
Consequently $M_5=M_T$.
The crossing $(5,8,9,+)$ would now force
$\alpha W=M_5u_8=M_TM_T^{-1}V=V$, impossible.

For $391$ put
\[
 (T,L)=(M_VW,JM_V^{-1}),\qquad
 (A,H)=(M_U^{-1}T,LM_U),\qquad t=F(T),\quad\tau=L(U).
\]
Then $A=T+(t/\mu)U$, $H=L-\tau F$.
The actual signed trace is
\[
 (9,6,5,-),(6,9,10,+),(4,10,11,-),
 (5,11,12,+),(11,7,8,+).
\]
The arc $5$ pair $D=M_W^{-1}(V,G)$ lies in the $V,W$ and $G,J$
planes. The arc $12$ pair is $M_D(A,H)$, so its $U,F$ coefficients
are $t/\mu,-\tau$. The crossing $(10,12,0,-)$ gives
\[
 u_{12}=\gamma(U-\tau T),\qquad
 f_{12}=\gamma^{-1}(F+(t/\mu)L).
\]
Hence $t=\mu\gamma$, $\tau=-\gamma^{-1}$ and $t\tau=-\mu$.
The displayed pair therefore equals $(A,H)$, so $M_{12}=M_A$.
The crossing $(12,8,9,-)$ and the arc $8$ pair $M_A(V,G)$ then
force $V=\beta W$, again impossible.

For $392$ define
\[
\begin{gathered}
 h=1+zZ/\mu,\quad t=-zx/\mu+hy,\quad\tau=ZX+hY,\\
 k=tX/\mu+z(h-d/\mu),\qquad K=-\tau x+(d+h)Z.
\end{gathered}
\]
The signed trace is
\[
 (12,6,5,+),(5,12,11,+),(4,11,10,-),
 (6,10,9,+),(10,7,8,-).
\]
Its paired coefficient vectors are
\[
\begin{array}{c|c|c}
i&u_i&f_i\\\hline
11&(0,-z/\mu,h)&(0,Z,h)\\
10&(t/\mu,-z/\mu,h)&(-\tau,Z,h)\\
9&(t/\mu,-z/\mu-k,h)&(-\tau,Z+K/\mu,h)\\
8&(Kt/\mu^2,1-Kz/\mu^2,Kh/\mu)&(k\tau,1-kZ,-kh).
\end{array}
\]
The crossing $(9,12,0,+)$, with
$p=f_9(W)$ and $q=J(u_9)$, gives
\[
 -pt/\mu=\gamma,\quad ph=1,\qquad
 -q\tau/\mu=\gamma^{-1},\quad qh/\mu=-1.
\]
Thus $p,q,t,\tau,h$ are nonzero. Its other two coefficients give
$k=-z/\mu$ and $K=-\mu Z$.
At $(11,8,9,+)$ the arc $11$ pair has zero $U,F$ coefficients.
The two corresponding coefficients therefore imply
\[
 Kt/\mu^2=\beta t/\mu,\qquad k\tau=-\beta^{-1}\tau.
\]
Only already nonzero $t,\tau$ are cancelled. We obtain
$K=\mu\beta$, $k=-\beta^{-1}$, hence
$zZ=kK=-\mu$ and $h=0$, contradicting $ph=1$.
The checker also verifies a constant-one integer Laurent identity
for these final necessary scalar equations.
\end{proof}

\section{The signed local folds}

For the four original diagrams $373,375,391,392$, the signed clasp
and its closing crossing have the same signs as in the companion
proof. It forces the two cross pairings with the third seed to vanish
and the remaining directed pairing product to be $-\mu$.
The resulting full matrices satisfy $M_4=M_0$ and $M_7=M_6$.
The changed seed lists and their signed determinant-one derivations
are retained explicitly in the input, including the changes needed
for $391$ and $392$.

For $3264,3266,3275,3276$, put
$(A,F)=(u_1,f_1)$, $(E,G)=(u_3,f_3)$, $(C,J)=(u_6,f_6)$ and
\[
 a=F(E),\quad b=G(A),\quad x=F(C),\quad X=J(A),\quad
 c=G(C),\quad C'=J(E),\qquad k=1+ab/\mu.
\]
The actual five local signs are
\[
 (3,1,0,+),(1,3,4,-),(0,3,2,+),(2,6,7,+),(7,1,2,+).
\]
They give
\[
 (u_2,f_2)=(-aA/\mu+kE,\ bF+kG).
\]
The root on arc $7$ has coefficient one on $C$.
The closing root relation $M_7A=\kappa u_2$ therefore first forces
$f_7(A)=0$, and then $A=\kappa u_2$.
Thus $k=0$ and $ab=-\mu$; moreover $\kappa=b\ne0$.
Under these relations the arc $7$ pair is
\[
                    (C-xA,\ J+(X/\mu)F).
\]
The vanishing evaluation is $X/\mu=0$, so $X=0$.
The closing covector equation is then $F+xJ=F$, forcing $x=0$.
Hence $M_7=M_6$. Directly,
\[
 (u_0,f_0)=(A-bE,F-b^{-1}G),\qquad
 (u_4,f_4)=(-b^{-1}u_0,-bf_0),
\]
so also $M_4=M_0$.
All inverses here are powers of $\mu$ or of a closing scalar already
proved nonzero.

For the eight A6 diagrams, put $(A,F)=(u_0,f_0)$,
$(E,G)=(u_1,f_1)$ and $(C,J)=(u_j,f_j)$, where $j=6$ or $8$.
Write $a=F(E)$, $b=G(A)$, $x=F(C)$, $X=J(A)$,
$c=G(C)$, $C'=J(E)$.
For the first knot in each pair below the local type is I; for the
second it is II. Every one of the five displayed exponents is negative:
\[
\begin{array}{c|l}
\mathrm I&(0,1,2,-),(j,2,3,-),(3,0,1,-),(1,3,4,-),(2,j,j',-)\\
\mathrm {II}&(j,1,2,-),(0,2,3,-),(3,0,1,-),(1,3,4,-),(2,j,j',-).
\end{array}
\]
In type I put $t=C'+aX/\mu$, $T=c-bx$.
The arc $3$ pair is
\[
 (E+(a/\mu)A+(t/\mu)C,\ G-bF-TJ).
\]
Let $H=\mu b-TX$, $L=(a+tx)/\mu$.
The closing relation $M_3^{-1}A=\kappa E$,
$FM_3=\kappa^{-1}G$ has coefficients
\[
 (1+aH/\mu^2,H/\mu,tH/\mu^2)=(0,\kappa,0),\quad
 (1+bL,-L,LT)=(0,\kappa^{-1},0).
\]
These force $t=T=0$, $a=-\mu\kappa^{-1}$, $b=\kappa$.
In type II put $t=a+C'x/\mu$, $T=b-cX$.
The same closing coefficient equations become
\[
 (1+tT/\mu,T,TC'/\mu)=(0,\kappa,0),\quad
 (1+tT/\mu,-t/\mu,tc/\mu)=(0,\kappa^{-1},0).
\]
They force $c=C'=0$, $a=t=-\mu\kappa^{-1}$, $b=T=\kappa$.
Both types give $ab=-\mu$, $M_4=M_0$ and $M_{j'}=M_j$.
No division by $1-\mu$ or $\mu^2-\mu+1$ occurs.

\section{Exact signed transfers and the upper theorem}

After removing the five local crossings and folding $0=4$ and $j=j'$,
the remaining eight signed relations match the earlier reduced systems
as follows. A sign $-$ means that every crossing exponent is reversed.
\[
\begin{array}{c|c|r}
\text{new knot}&\text{old reduced system}&\text{global sign}\\\hline
3264&373&-\\3266&375&-\\3275&391&-\\3276&392&-\\
304&375&-\\3232&375&+\\362&375&+\\3259&375&-\\
623&373&-\\3374&373&+\\629&373&-\\3375&373&+
\end{array}
\]
All indices denote $13a$ knots. The first four rows use the identity
relabeling. The remaining relabelings are explicit bijections in the
input; the checker reconstructs and compares the signed crossing sets.
It permits reversal of an under-arc pair only with reversal of its
exponent. It does not reuse an unsigned incidence match as a signed one.

Global sign reversal causes no loss: replace every reduced matrix by
its inverse and replace the parameter by $\mu^{-1}$.
Relation by relation, $N_b=N_o^{-\varepsilon}N_aN_o^{\varepsilon}$.
The root directions are unchanged, and all covectors are multiplied
by the nonzero scalar $-\mu^{-1}$. Thus both spans and their
independence are preserved. The pair evaluation becomes
$1-\mu^{-1}$, precisely the same model at the reciprocal parameter.

The required independence follows before the reduced equations are
used. In the first four cases, the changed four-seed root basis gives
$U=A-B,V=C,W=D$; the covector basis gives $F_A-F_B,F_C,F_D$.
Both triples are independent. For the alternate four cases the triples
are $A-bE,C,Q$ and $F-b^{-1}G,J,F_Q$, subsets of triangular changes
of the two four-seed bases; here $Q$ is the last retained changed seed.
For the eight A6 cases, the reduced triples are $u_0,u_6,u_9$;
$u_0,u_6,u_{11}$; or $u_0,u_5,u_8$, with the corresponding covectors,
as specified by their changed seed lists. Each is a subset of its
changed basis. The local folds identify other meridian matrices and
do not remove these independent triples.
The local folds and signed transfers therefore contradict
Lemma~\ref{lem:cores}. This proves Theorem~\ref{thm:all16}.
The proof permits singular proper pairing submatrices, $\mu=1$ and
primitive sixth roots of unity. Only matrix invertibility and closing
scalars already proved nonzero are used for cancellation.

For a representation $\rho:G_K\to GL(V)$, where
$G_K=\pi_1(S^3\setminus K)$, its displacement span means the vector
subspace
\[
              \sum_{g\in G_K}\operatorname{im}(\rho(g)-I)\subset V.
\]
It is a span of displacement images, rather than a space of matrices.

\begin{corollary}\label{cor:semisimple}
For every $K\in\mathcal K_{16}$, any finite-dimensional completely
reducible representation over any field with meridian displacement
rank one has common-fixed-space codimension and displacement-span
dimension at most three. This applies in particular to every such
finite-image representation whose image order is nonzero in the field,
including every finite-image complex representation.
\end{corollary}
\begin{proof}
Decompose into irreducible summands. The ranks of the meridian
displacements add, so exactly one summand is active. On every other
summand the meridian is the identity, and normal generation by a
meridian makes the whole knot group act trivially there.
The active irreducible summand has dimension at most three by
Theorem~\ref{thm:all16}. All displacements lie in that summand,
which bounds their span, and the inactive summands are pointwise fixed,
which bounds common-fixed-space codimension.

For the finite-image assertion, write the image as $H$ and assume
$|H|$ is invertible in the field. If $W$ is invariant, choose any
projection $P$ onto $W$. The average
\[
                 \overline P=|H|^{-1}\sum_{h\in H}hPh^{-1}
\]
is an equivariant projection onto $W$; its kernel is an invariant
complement. Induction on dimension gives complete reducibility.
\end{proof}

\begin{corollary}\label{cor:augmentation}
For every $K\in\mathcal K_{16}$, every field-valued $U=1$ cord
augmentation with $\epsilon(\mu)\ne1$ has rank at most three, with
arbitrary nonzero longitude value.
\end{corollary}
\begin{proof}
Apply Cornwell's arbitrary-field augmentation--irreducible-KCH
correspondence, meridian/longitude identification and rank convention
\cite[Theorem 1.2, Definition 1.8, Proposition 2.12 and Lemma 3.14]{cor17}.
The associated irreducible representation has dimension at most three
by Theorem~\ref{thm:all16}.
Positive augmentation witnesses and meridional-rank lower bounds are
separate certificates and are not premises of this augmentation upper bound.
\end{proof}

\begin{corollary}\label{cor:gap}
For every $K\in\mathcal K_{16}$,
\[
              \operatorname{ar}(K,\mathbb C)=3<4=\mr(K)=b(K).
\]
\end{corollary}
\begin{proof}
The old eight rank-three reflective witnesses are retained in
\path{REFLECTIVE_AUGMENTATION_RANK_GAP.tex}, Theorem~3; the other
eight lower witnesses are Theorem~1, \texttt{thm:attainment}, of
\path{A6_RANK_THREE_AUGMENTATIONS.tex}. They give $U=1$ cord
augmentations of rank three at the common parameter $\mu=-1$.
Together with Corollary~\ref{cor:augmentation} they prove augmentation
rank exactly three. The independent seven-sheet and A6 affine
certificates supply the meridional-rank and bridge equalities.
Only the positive attainment theorem from the A6 lower package is
used here, not its corollary importing the present upper bound.
\end{proof}

\begin{corollary}\label{cor:sums}
For every $K\in\mathcal K_{16}$ and every integer $s\ge0$, let
$K_s=K\mathbin{\#}(\mathbin{\#}^{s}3_1)$. Then
\[
          \operatorname{ar}(K_s,\mathbb C)=s+3<s+4=\mr(K_s)=b(K_s).
\]
\end{corollary}
\begin{proof}
Each seed has a rank-three augmentation at $\mu=-1$, and the trefoil
has a rank-two augmentation at that same parameter, as explicitly
verified in the companion $373$ package. Cornwell's Theorem~1.9 gives
the lower construction and the matching upper additivity of
augmentation rank, hence $\operatorname{ar}(K_s,\mathbb C)=s+3$.
The separate normalized-cover consequence
\path{INVOLUTION_COVER_SUM_FORMULAS.tex},
Corollary~\texttt{cor:hundredone}, gives $\mr(K_s)=b(K_s)=s+4$ for
these sixteen seeds. No unrestricted meridional-rank additivity is
assumed.
\end{proof}

Strict augmentation-rank gaps were already known: Cornwell's
Corollary~1.12 gives $2=\operatorname{ar}(K,\mathbb C)<\mr(K)=3$
for $8_{16},8_{17},10_{91},10_{94}$. No first-counterexample or
worldwide-priority claim is made here.

\section{Exact replay and limits}

Run \path{verify_generic_kch_extensions.py}; its default mode uses
only the standard library. The input
\path{generic_kch_extensions_inputs.json} retains all sixteen literal
PDs, original and changed seed systems, signed crossings, signed basis
words, census bindings and every reduced-system relabeling.
The checker imports only pinned exact arithmetic/diagram helpers and
the first reduced contradiction from the frozen $373$ package.
It verifies the three additional free-directed-pairing traces, the
raw local folds, both basis determinants, the signed transfers and an
additional integer Laurent identity with constant one.
The optional \texttt{--census} mode rereads every named PD and DT from
the pinned independently packaged census. Corrupted signs and named
maps are rejected. No finite parameter sampling is used.

The upper theorem uses irreducibility, and Corollary~\ref{cor:semisimple}
extends its codimension consequence to completely reducible representations.
The two independent families remain essential to the main proof; no
fixed-space codimension bound is asserted for arbitrary nonsemisimple
rank-one representations over general fields.
The augmentation statement uses the exact published correspondence and
does not address augmentations at $\epsilon(\mu)=1$ outside its rank
construction. All algebra and diagram checks are internal; formal
verification and external specialist review are not claimed.

\begin{thebibliography}{9}
\bibitem{cor17} C. R. Cornwell, KCH representations, augmentations,
and $A$-polynomials, J. Symplectic Geom. 15 (2017), 983--1017.
Retained author version arXiv:1310.7526v3, Theorem 1.2,
Definition 1.8, Proposition 2.12, Lemma 3.14, Theorem 1.9, Corollary 1.12.
\url{https://arxiv.org/abs/1310.7526v3}.
\end{thebibliography}
\end{document}
